\documentclass[12pt]{article}
\usepackage{amsmath, amssymb, bm}
\usepackage{geometry}
\usepackage{hyperref}
\geometry{margin=1in}

\begin{document}

\title{BHRuler Equations}
\author{Stephen L. Nagy \\ \small{Independent Researcher}}
\date{March 2026 \\ Version 2.1.1-survey}
\maketitle

\section*{Introduction}

This note lists the equations currently implemented by \textbf{BHRuler} (the single-file
\texttt{bhruler.py}) and maps them to the output columns written by the code.

This document is intentionally limited to \emph{implemented} framework math:
\begin{itemize}
    \item gravitational scales and the Schwarzschild ISCO invariant,
    \item Kerr spin corrections via Bardeen--Press--Teukolsky,
    \item dual accretion prescriptions ($\eta$-bridge and ADAF/RIAF),
    \item horizon magnetic field and Blandford--Znajek jet-power scaling,
    \item environment metrics,
    \item TDE fallback and the logistic capture boundary.
\end{itemize}

Paper-level discussion that is not yet a first-class code output (for example, the full uncertainty-propagation machinery) is not treated here as implemented CLI math.

\section*{Notation and Inputs}

\begin{itemize}
    \item $M$: black-hole mass
    \item $a_* \in [-1,1]$: dimensionless spin parameter
    \item $\lambda_{\rm Edd} \equiv L_{\rm bol}/L_{\rm Edd}$: Eddington ratio
    \item $\eta_{\rm eff}$: effective radiative efficiency
    \item $\kappa$: ADAF/RIAF proportionality constant defined by
    \[
    \lambda_{\rm Edd} = \kappa \dot m^2
    \]
    \item $\dot m$: dimensionless accretion rate
    \item $\sigma$: stellar velocity dispersion
    \item $R_e$: galaxy effective radius
    \item $R_*, M_*$: stellar radius and stellar mass for the fallback-time scaling
\end{itemize}

\section*{Physical Constants (SI)}

\[
G = 6.67430\times10^{-11}\ {\rm m^3\,kg^{-1}\,s^{-2}}
\]
\[
c = 2.99792458\times10^{8}\ {\rm m\,s^{-1}}
\]
\[
M_\odot = 1.98847\times10^{30}\ {\rm kg}
\]
\[
1\,{\rm pc} = 3.085677581491367\times10^{16}\ {\rm m}
\]

\section{Core Gravitational Scalings}

\subsection*{Gravitational time and length}
\[
t_g = \frac{GM}{c^3},
\qquad
r_g = \frac{GM}{c^2},
\qquad
r_s = 2r_g = \frac{2GM}{c^2}.
\]

\subsection*{Schwarzschild ISCO baseline}
\[
r_{\rm ISCO}^{\rm Schw} = 6r_g
\]
\[
f_{\rm ISCO}^{\rm Schw} = \frac{c^3}{6^{3/2}\,2\pi GM}.
\]

\subsection*{Mass-invariant sanity check}
BHRuler evaluates the dimensionless invariant
\[
f_{\rm ISCO}^{\rm Schw} t_g
= \frac{c^3}{6^{3/2}\,2\pi GM}\cdot\frac{GM}{c^3}
= \frac{1}{6^{3/2}\,2\pi}
\approx 0.010829122239.
\]

This quantity should be independent of mass and is used as a built-in validation check.

\section{Kerr Spin Corrections (Bardeen--Press--Teukolsky)}

\subsection*{ISCO radius in units of $r_g$}
Define
\[
Z_1 = 1 + (1-a_*^2)^{1/3}
\left[(1+a_*)^{1/3} + (1-a_*)^{1/3}\right],
\]
\[
Z_2 = \sqrt{3a_*^2 + Z_1^2}.
\]

Then
\[
r_{\rm ISCO}(a_*) =
\begin{cases}
3 + Z_2 - \sqrt{(3-Z_1)(3+Z_1+2Z_2)}, & \text{prograde} \\[6pt]
3 + Z_2 + \sqrt{(3-Z_1)(3+Z_1+2Z_2)}, & \text{retrograde}
\end{cases}
\]
where $r_{\rm ISCO}$ is expressed in units of $r_g=GM/c^2$.

\subsection*{Kerr ISCO orbital frequency}
For the chosen orbital branch,
\[
f_{\rm ISCO}^{\rm Kerr}(a_*) =
\frac{c^3}{2\pi GM}
\frac{1}{r_{\rm ISCO}^{3/2}(a_*) \pm a_*},
\]
with
\[
+\ a_* \quad \text{for prograde},
\qquad
-\ a_* \quad \text{for retrograde}.
\]

\subsection*{Spin diagnostic ratio}
BHRuler reports
\[
f_{\rm ratio}
=
\frac{f_{\rm ISCO}^{\rm Kerr}}{f_{\rm ISCO}^{\rm Schw}}.
\]

The CLI computes the prograde branch by default for its standard output columns.

\subsection*{ISCO orbital period}
When survey-style output aliases are written, the code also reports
\[
P_{\rm orb,ISCO} = \frac{1}{f_{\rm ISCO}^{\rm Kerr}}.
\]

\section{Accretion Prescriptions}

BHRuler supports four branch modes:
\begin{itemize}
    \item \texttt{core}: geometry, timing, spin, environment, and TDE only
    \item \texttt{eta\_bridge}: efficient/thin-disk-like accretion
    \item \texttt{adaf}: radiatively inefficient ADAF/RIAF-like prescription
    \item \texttt{both}: compute both accretion branches
\end{itemize}

\subsection*{(a) Efficiency bridge}
The accretion rate is defined by
\[
\dot m = \frac{\lambda_{\rm Edd}}{\eta_{\rm eff}}.
\]

The horizon magnetic field is modeled as
\[
B_H^{(\eta)} \simeq 4\times10^4\ {\rm G}\;
\dot m^{1/2}
\left(\frac{10^9 M_\odot}{M}\right)^{1/2}.
\]

Equivalently,
\[
B_H^{(\eta)} \simeq 4\times10^4\ {\rm G}\;
\left(\frac{\lambda_{\rm Edd}}{\eta_{\rm eff}}\right)^{1/2}
\left(\frac{10^9 M_\odot}{M}\right)^{1/2}.
\]

\subsection*{(b) ADAF/RIAF}
The ADAF closure is
\[
\lambda_{\rm Edd} = \kappa \dot m^2
\quad\Longrightarrow\quad
\dot m = \sqrt{\frac{\lambda_{\rm Edd}}{\kappa}}.
\]

BHRuler uses the same field scaling in terms of $\dot m$,
\[
B_H^{(\rm ADAF)} \simeq 4\times10^4\ {\rm G}\;
\dot m^{1/2}
\left(\frac{10^9 M_\odot}{M}\right)^{1/2},
\]
hence
\[
B_H^{(\rm ADAF)} \simeq 4\times10^4\ {\rm G}\;
\left(\frac{\lambda_{\rm Edd}}{\kappa}\right)^{1/4}
\left(\frac{10^9 M_\odot}{M}\right)^{1/2}.
\]

\subsection*{Interpretive note}
The leading physical dependence in both branches is
\[
B_H \propto \dot m^{1/2} M^{-1/2}.
\]
The numerical prefactor is treated as MAD-motivated and order-of-magnitude.

\section{Blandford--Znajek Jet Power}

Given a horizon-scale field estimate $B_H$, BHRuler reports the order-of-magnitude jet scaling
\[
P_{\rm BZ} \approx 10^{45}\ {\rm erg\,s^{-1}}
\left(\frac{a_*}{0.9}\right)^2
\left(\frac{B_H}{10^4\,{\rm G}}\right)^2
\left(\frac{M}{10^9M_\odot}\right)^2.
\]

This is a phenomenological normalization anchored to the usual
\[
P_{\rm BZ} \propto a_*^2 \Phi_{\rm BH}^2
\]
scaling, with the BHRuler normalization tied to MAD-motivated GRMHD expectations.

\section{Environment and Host Coupling}

\subsection*{Sphere of influence}
\[
r_{\rm infl} = \frac{GM}{\sigma^2}.
\]

In the code this is returned in parsecs.

\subsection*{Host-normalized influence radius}
If an effective radius is supplied,
\[
\frac{r_{\rm infl}}{R_e}
\]
is also reported, with the code converting parsecs to kiloparsecs before division.

\section{Tidal-Disruption Scalings}

\subsection*{Fallback time}
For a solar-type normalization,
\[
t_{\rm fb} \approx 41~{\rm d}
\left(\frac{M}{10^6M_\odot}\right)^{1/2}
\left(\frac{R_*}{R_\odot}\right)^{3/2}
\left(\frac{M_*}{M_\odot}\right)^{-1}.
\]

The implemented default is $R_*=R_\odot$ and $M_*=M_\odot$.

\subsection*{Logistic TDE capture boundary}
BHRuler defines a logistic capture factor
\[
S(M;a_*) =
\frac{1}{1+\exp(-x)},
\]
with
\[
x = \frac{\log_{10} M - \log_{10} M_{\rm crit}(a_*)}{w}.
\]

The spin-shifted midpoint is
\[
M_{\rm crit}(a_*) = M_{\rm crit,0}\,(1-0.6\,a_*),
\]
where the code clamps the spin shift to the interval $a_* \in [0,1]$ for this purpose.

Default values:
\[
M_{\rm crit,0} = 3\times10^7\,M_\odot,
\qquad
w = 0.15\ {\rm dex}.
\]

The interpretation is:
\[
S \rightarrow 0 \quad \text{below the midpoint (disruption favored)},
\]
\[
S \rightarrow 1 \quad \text{above the midpoint (direct capture favored)}.
\]

The boolean output is
\[
\texttt{tde\_possible} = (S < 0.5).
\]

\subsection*{Heuristic status}
The factor $(1-0.6a_*)$ is used as a tunable heuristic rather than a precision GR fit.

\section{Implemented Output Mapping}

\subsection*{Core outputs}
\begin{itemize}
    \item \texttt{t\_g\_s} $\leftrightarrow t_g$
    \item \texttt{r\_s\_km} $\leftrightarrow r_s$
    \item \texttt{f\_ISCO\_Schw\_Hz} $\leftrightarrow f_{\rm ISCO}^{\rm Schw}$
    \item \texttt{fISCO\_tg\_invariant} $\leftrightarrow f_{\rm ISCO}^{\rm Schw} t_g$
\end{itemize}

\subsection*{Spin-aware outputs}
\begin{itemize}
    \item \texttt{r\_ISCO\_rg} $\leftrightarrow r_{\rm ISCO}(a_*)/r_g$
    \item \texttt{f\_ISCO\_Kerr\_Hz} $\leftrightarrow f_{\rm ISCO}^{\rm Kerr}(a_*)$
    \item \texttt{f\_ratio} $\leftrightarrow f_{\rm ISCO}^{\rm Kerr}/f_{\rm ISCO}^{\rm Schw}$
\end{itemize}

\subsection*{Environment outputs}
\begin{itemize}
    \item \texttt{rinfl\_pc} $\leftrightarrow r_{\rm infl}$
    \item \texttt{rinfl\_over\_Re} $\leftrightarrow r_{\rm infl}/R_e$
    \item \texttt{t\_fb\_days} $\leftrightarrow t_{\rm fb}$
\end{itemize}

\subsection*{TDE output}
\begin{itemize}
    \item \texttt{tde\_possible} $\leftrightarrow (S<0.5)$
\end{itemize}

\subsection*{Accretion and jet outputs}
\begin{itemize}
    \item \texttt{B\_H\_G\_eta}, \texttt{P\_BZ\_erg\_s\_eta}
    \item \texttt{B\_H\_G\_adaf}, \texttt{P\_BZ\_erg\_s\_adaf}
\end{itemize}

\subsection*{Warning outputs}
If a requested accretion branch is missing its required parameter, the code writes
\begin{itemize}
    \item \texttt{warn\_eta\_bridge}
    \item \texttt{warn\_adaf}
\end{itemize}
and fills the branch-specific outputs with NaNs.

\subsection*{Survey-compatible aliases}
When survey-like input is detected, the code may also emit alias columns such as
\begin{itemize}
    \item \texttt{Name}
    \item \texttt{Mass\_Msun}
    \item \texttt{Spin\_a}
    \item \texttt{Grav\_Time\_tg\_s}
    \item \texttt{Schwarzschild\_Rad\_km}
    \item \texttt{Orbital\_Period\_s}
    \item \texttt{Invariant\_fISCO\_tg}
\end{itemize}
for easy round-trip comparison with the Survey V2 table.

\section{Validation Checks}

\begin{itemize}
    \item \textbf{Invariant check:}
    \[
    f_{\rm ISCO}^{\rm Schw} t_g = \frac{1}{6^{3/2}2\pi}
    \]
    should remain constant across mass.
    
    \item \textbf{Kerr branch check:}
    prograde spin drives $r_{\rm ISCO}$ downward and raises $f_{\rm ISCO}^{\rm Kerr}$.
    
    \item \textbf{Unit check:}
    \[
    r_s = 2r_g,
    \qquad
    r_{\rm infl}\propto \frac{M}{\sigma^2}.
    \]
    
    \item \textbf{Branch logic check:}
    \texttt{core} should produce geometry/spin/environment/TDE outputs without requiring
    $\lambda_{\rm Edd}$, $\eta_{\rm eff}$, or $\kappa$.
\end{itemize}

\section*{References (Core Methods)}

\begin{itemize}
    \item Bardeen, Press \& Teukolsky (1972) --- Kerr ISCO and equatorial geodesics
    \item Blandford \& Znajek (1977) --- electromagnetic extraction of black-hole spin energy
    \item Shakura \& Sunyaev (1973) --- efficient thin-disk phenomenology
    \item Narayan \& Yi (1994, 1995) --- ADAF/RIAF scaling
    \item Tchekhovskoy et al. (2011) --- MAD-saturated GRMHD jet normalization
    \item McKinney et al. (2012) --- GRMHD jet and magnetic-flux calibration context
    \item Kesden (2012) --- spin dependence of tidal disruption thresholds
    \item Servin \& Kesden (2017) --- relativistic disruption/capture refinements
    \item Rees (1988) and Phinney (1989) --- canonical TDE fallback-time scaling
\end{itemize}

\end{document}
